% Section 10 — Conclusion
\section{Conclusion: the search for Move 42}
\label{sec:conclusion}

Move~37 showed that a machine could make a move no person would have made and win with it. It took
a search to find it. InnovationZero asks whether the search can be learned and then removed: whether
a model can read a dataset and write the algorithm for it, as working code, before trying anything.

We have made that question precise and testable. We defined zero by counting --- one committed
program, no executions, no revisions --- and separated three tracks that are usually confused. We gave
frontier models a closed-world challenge and found that they can invent working algorithms by
reasoning alone, that the best of them is very good, and that all of them reinvent the same
textbook design while keeping a small fraction of the tokens they process. We showed that a
data-conditioned engine already writes a different, standalone program for every dataset, about as
good as the best frontier invention on our pilot panel, and that strong programs for these datasets
exist in the space it searches. And we proved that the most obvious learned proposer, one that
chooses from a catalog, can never answer the question, which tells us what the one that can must
look like.

The same construction is not specific to regression. Any problem whose solution is a program and
whose outcome can be measured admits the same contract, the same tracks and the same ladder: binary
and multiclass classification, scored by held-out discrimination and calibration; forecasting,
scored out of time; combinatorial optimization, scored by objective value under a time limit; and,
because a system's own proposal and selection procedures are programs with measurable outcomes, the
improvement of the system itself --- the last rung of the ladder, and the one that should be
claimed last. The repository already contains classification heads and metrics and time-series
fixtures; none of these extensions has yet been qualified, and we list them as directions, not
results.

We call the class of models that write algorithms from data and learn from what those algorithms do
\emph{algorithmic foundation models}. Language models learned to write about solutions. The open
question is whether a model can learn, from the outcomes of experiments, to write the solution itself
--- and the first program it writes that works, that no archive contained and that no frontier model
would have proposed, will be its Move~42. We have released the challenge, the entrants, the pilots and
the protocol so that anyone can look for it.

\section*{Acknowledgements and disclosures}
\addcontentsline{toc}{section}{Acknowledgements and disclosures}
The frontier-model entrants were written by the models named for them, in sessions run by the
repository's owner; integration and verification were carried out by coding agents working in the
repository, whose receipts are cited. This paper was drafted with the assistance of AI models,
and every number in it was checked against the repository files listed in \cref{app:claims}. Jepson
Taylor is the founder of Veox AI, which develops the system described. The kernel language and the
trained proposer are proprietary; \cref{app:disclosure} lists what is and is not released.
